% Pursuing Stacks (cotes 134-2 to 134-8): archive cote 134-5, site page 82 (PDF page 83), typescript page 269.
% typescript. Draft compared: ps4.tex:4190-4271.
% AI-transcribed by Claude Opus 5.5 (claude-opus-5-5, adaptive thinking, effort high), Message Batches API, batch ps2 (msgbatch_01NVwr1Yuo8qJf5rykeGBXPK), 2026-11-12; not yet checked by eye.
\scanpage{134-5/82}
\mspage{269}

\uline{Application to relation between contractibility and objects} \uline{Hom}(X,Y).

I would like to review here a few things along these lines, which were
somewhat scattered in the notes before (section 51 E) p.121 and section 57
p. 143 notably), and at times came out awkwardly because of unadequate
conceptual background. We assume M endowed with a contractibility structure
$\mathrm{M}_c$ , and use the notations of the previous section\unsure{;|,} especially concerning
the subsets of Ob M and of \unsure{Fl}(M)
\[ (1) \qquad \mathrm{M}_c \hand{\subset} \mathrm{M}_\infty \hand{\subset} \mathrm{M}_{\underline{W}} \hand{\subset} \mathrm{M}_o \ , \quad \mathrm{W}_c \hand{\subset} \mathrm{W}_\infty \hand{\subset} \underline{\mathrm{W}}_{\mathrm{M}} \hand{\subset} \xout{x}\, \mathrm{W}_o \quad . \]
Let X be an object of M, such that the object \xout{Hom}
\[ (2) \qquad \mathrm{I} = \underline{\mathrm{Hom}}(\mathrm{X},\mathrm{X}) \]
exists in M (NB a priori it is an object in M\^{}, we assume it to be repre-%
sentable). We assume that X is endowed with a section c , hence a section
\xout{of} $\hand{\delta}_1$ of I , corresponding to the constant endomorphism of X with value c,
We'll denote by $\hand{\delta}_o$ the section of I corresponding to the identity of X.
Thus
\[ (3) \qquad \underline{\mathrm{I}} = (\mathrm{I},\hand{\delta}_o,\hand{\delta}_1) \]
becomes an interval of M, and the composition law of I = \uline{Hom}(X,X) turns
it into a \uline{multiplicative} \uline{interval} , admitting respectively $\hand{\pi}_o$ and $\hand{\pi}_1$ as
left unit and left zero element. Then the following conditions are equi-%
valent:

(\xout{xx}i) X is contractible.

(ii) \uline{Hom}(X,X) is contractible, \hand{i.e. $\mathrm{I}\in \mathrm{M}_c$ .}

(ii') \uline{Hom}(X,X) is 0-connected, \hand{i.e. $\mathrm{I}\in \mathrm{M}_o$ .}

(ii") \xout{Hom} The two sections $\hand{\delta}_o,\hand{\delta}_1$ of $\underline{\mathrm{I}}$=\uline{Hom}(X,X) are homotopic.

(iii) For any object Y in M such that \uline{Hom}(Y,X) exists in M, \uline{Hom}(Y,X)
is contractible.

(iii') For any Y as in (iii), \uline{Hom}(Y,X) is \inserted{0-}connected.

(iv) For any Y in M such that \uline{Hom}(X,Y) exists in M, the canonical
map
\[ (4) \qquad \mathrm{Y} \hand{\longrightarrow} \underline{\mathrm{Hom}}(\mathrm{X},\mathrm{Y}) \]
is a homotopism (more accurately still, Y as a subobject of
\uline{Hom}(X,Y) is a deformation ret\xout{ra}act). Moreover, X is 0-connected.

(iv') For any Y in M as in (iv), the map (4) induces a bijection \xout{xx}
\[ \xout{x} \qquad \hand{\pi}_o(\mathrm{Y}) \hand{\to} \hand{\pi}_o(\underline{\mathrm{Hom}}(\mathrm{X},\mathrm{Y})) \quad , \]
moreover X is 0-connected.

NB Of course, the 0-connectedness notion and the functor $\hand{\Pi}_o$ used \xout{her}

\margin{82}
